\section{考察と限界}
\label{sec:discussion}

\paragraph{アノテーション削減の源泉。}
本システムは幾何を不要にするのではなく、そのコストの償却方法を変える。手作業では、画像ごとに
キャリブレーションするか、画像上の点をクリックして指定する。提案する作業手順では、動画取得、再構成、
および一度の監査済みシーンアライメントにコストをかける。その後、カメラ配置、姿勢、内部パラメータ、
点の同一性、ライン／セグメンテーション教師信号、深度、幾何学的妥当性、およびレンダラ支持を、
受理されたレンダリングごとに計算する。削減効果は、レンダリングしたビューの数と多様性に応じて増大する。
これにより、局所ショートカットを打破するうえで特に有用な、コートの一部しか含まないビュー、
高所からのビュー、光軸を外したビュー、および複数コートを含むビューを低コストで追加できる。

\paragraph{3DGSは必要だが十分ではない。}
新規視点レンダリングだけでは、メートル尺度もコートの意味構造も得られない。見栄えが良くても、
アライメントされていないGaussianシーンは、並進や物理キーポイントの同一性に対する教師信号を
与えられない。逆に、レンダラを伴わないコート変換が生成するのはラベルだけであり、新規姿勢における
現実的な画素は生成しない。本研究の貢献は、単一の変換グラフを介して両者を統合し、尺度、トポロジー、
またはホールドアウトデータから得たエビデンスが弱い場合には、その統合結果を棄却することにある。

\paragraph{アライメントの初期化と再構成の失敗。}
アライメントは学習済みラインマップを用いるため、その誤りモードには既存のライン検出器が持つ誤りも
含まれる。疎な再構成点では地面を十分に支持できない場合があり、反射するラインや遮蔽されたラインを
見落とす可能性もある。また、反復して並ぶ隣接コートは、向きが90度ずれる誤りやコートの同一性を
取り違える誤りを引き起こし得る。共通尺度、テンプレート全体、トポロジー、およびホールドアウトの
各ゲートは、これらの誤りの伝播を抑える一方で、利用できそうに見えるシーンも棄却する。これらは
品質管理であって、受理されたすべての変換が正確であることの証明ではない。

\paragraph{レンダリングのドメインギャップ。}
3DGS再構成は撮影施設に近い外観を保つが、浮遊物、ぼけ、一時的な物体の欠落、または撮影カメラ群が
覆う範囲を大きく外れた領域における信頼性の低い外挿を含み得る。そのため、軌道半径は撮影時の
カメラオフセットを基準に定め、レンダリング前に幾何条件を検査する。学習時の測光的データ拡張により
さらに変動を与えるが、あらゆる放送用カメラ、コート表面、気象条件、または圧縮パイプラインに対する
ドメインギャップを解消すると本稿は主張しない。

\paragraph{カメラモデル。}
姿勢ヘッドは二つの焦点距離が等しいことを仮定し、データ拡張から与えられる主点を用いる。
歪み、スキュー、または移動する主点は予測しない。ラベルの一意性を維持するため、未対応の変換は
棄却するが、より広範なカメラモデルには拡張したヘッドと投影損失が必要である。
カメラ側エンドの正準化でも、反射を含まない二つの回転 \(I\) と \(R_z(\pi)\) のいずれかを仮定する。
中央平面付近の著しく曖昧なビューは、暗黙に割り当てず除外しなければならない。

\paragraph{ショートカット耐性の範囲。}
大域クエリと再投影目的関数は解の集合を制約するが、ニューラルネットワークがあらゆる非本質的な
手がかりを無視することを証明するものではない。ネットワークは、姿勢と相関するレンダリング由来の人工的特徴や
施設固有のテクスチャを利用する可能性があり、二つの分岐が共適応する可能性もある。意図した機構を
検証するには、正解に対する直接損失、独立した姿勢再投影指標、勾配停止アブレーション、最悪層の報告、
および局所文脈のみを残す診断が必要である。キャリブレーション済みのホールドアウト実データによる評価が、
依然として最も重要な未実施項目である。

\paragraph{計算コストとデータガバナンス。}
再構成と新規視点レンダリングは学習前のコストを増加させるが、一つのシーンから多数のタスク用
データセットとカメラビューを生成することで、そのコストを償却できる。素材動画と再構成された外観には、
人物や施設を特定し得る情報が含まれる可能性がある。データセットを公開する際には、ソフトウェアの
オープンソースライセンスとは独立に、素材のライセンスとプライバシー制約を遵守しなければならない。
